\chapter{Differentiation}

\emph{Differentiation} is a core concept in calculus. It describes the rate at which a quantity changes 
and provides tools for analyzing and modeling change in real-world phenomena.

\section{Definition of the Derivative}

The \emph{derivative} of a function \(f\) at a point \(x\) is defined as the limit:

\[
    f'(x) = \lim_{h \to 0} \frac{f(x + h) - f(x)}{h}
\]

or

\[
    f'(x) = \lim_{h \to a}\frac{f(x) - f(a)}{x - a}
\]

provided this limit exists.

Geometrically, \(f'(a)\) is the slope of the tangent line to the graph of \(f\) at \(x = a\).

\section{Derivative Rules}

\emph{Power Rule:}

\[
    \frac{d}{dx} x^n = nx^{n - 1}, \quad \text{for } n \in \R
\]

\emph{Constant Rule:}

\[
    \frac{d}{dx} c = 0, \quad \text{for constant } c
\]

\emph{Constant Multiple Rule:}

\[
    \frac{d}{dx} [c \cdot f(x)] = c \cdot f'(x)
\]

\emph{Sum and Difference Rule:}

\[
    \frac{d}{dx} [f(x) \pm g(x)] = f'(x) \pm g'(x)
\]

\emph{Product Rule:}

\[
    \frac{d}{dx} [f(x)g(x)] = f'(x)g(x) + f(x)g'(x)
\]

\emph{Quotient Rule:}

\[
    \frac{d}{dx} \left( \frac{f(x)}{g(x)} \right) = \frac{f'(x)g(x) - f(x)g'(x)}{{g(x)}^2}
\]

\emph{Chain Rule:}

\[
    \frac{d}{dx} f(g(x)) = f'(g(x)) \cdot g'(x)
\]

\section{Trigonometric Derivatives}

\begin{align*}
    \frac{d}{dx} \sin(x) &= \cos (x) \\
    \frac{d}{dx} \cos(x) &= -\sin(x) \\
    \frac{d}{dx} \tan(x) &= \sec^2(x) \\
    \frac{d}{dx} \cot(x) &= -\csc^2(x) \\
    \frac{d}{dx} \sec(x) &=  \sec(x)\tan(x) \\
    \frac{d}{dx} \csc(x) &= -\csc(x) \cot(x)
\end{align*}

\section{Derivatives of Inverse Trigonometric Functions}

\begin{align*}
    \frac{d}{dx} \arcsin(x) &= \frac{1}{\sqrt{1 - x^2}} \\
    \frac{d}{dx} \arccos(x) &= -\frac{1}{\sqrt{1 - x^2}} \\
    \frac{d}{dx} \arctan (x) &= \frac{1}{1 + x^2} \\
    \frac{d}{dx} \mathrm{arcot}(x) &= -\frac{1}{1 + x^2} \\
    \frac{d}{dx} \mathrm{arcsec}(x) &= \frac{1}{|x|\sqrt{x^2 - 1}} \\
    \frac{d}{dx} \mathrm{arccsc}(x) &= -\frac{1}{|x|\sqrt{x^2 - 1}}
\end{align*}

\section{The Differential}

The \emph{differential} \(dy\) of a function \(y = f(x)\) is defined as:

\[
    dy = f'(x) \, dx.
\]

This linear approximation estimates the change in \(y\) for a small change in \(x\).

\section{The Secant Equation and Graphical Meaning}

The \emph{secant line} through the points \((x_0, f(x_0))\) and \((x_0 + h, f(x_0 + h))\) has slope

\[
    \frac{f(x_0 + h) - f(x_0)}{h}.
\]

As \(h \to 0\), the secant line approaches the tangent line. Graphically, this means

\[
    \lim_{h \to 0} \text{slope of secant} = \text{slope of tangent} = f'(x) = \Delta_h.
\]

The secant line equation is given by

\[
    L(x) = f(x) + \Delta_h f(x_0)(x - x_0).
\]

\section{Optimization via the Derivative}

To find local or global extrema:

\begin{enumerate}

    \item Compute \(f'(x)\).

    \item Solve \(f'(x) = 0\) to find the critical points.

    \item Use the first or second derivative test to classify them.

    \item Evaluate the endpoints if optimizing over a closed interval.

\end{enumerate}

\textbf{Example I:}

Maximize the area \(A = x(10 - 2x)\):

\[
    A = 10x - 2x^2, \quad A' = 10 - 4x, \quad A' = 0 \Rightarrow x = 2.5.
\]

\textbf{Example II:}

Given a number \( c > 0 \), we want to factor it into two positive numbers \( x \) and \( y \) such that 
their product is \( c \) and their sum \( S = x + y \) is minimal.

Since \( x \cdot y = c \), we can express one variable in terms of the other, e.g.,

\[
    y = \frac{c}{x}.
\]

Substituting into the sum,

\[
    S(x) = x + \frac{c}{x}.
\]

To minimize \( S(x) \), we take the derivative:

\[
    S'(x) = 1 - \frac{c}{x^2}.
\]

Setting the derivative to zero gives

\[
    1 - \frac{c}{x^2} = 0 \quad \Rightarrow \quad x^2 = c \quad \Rightarrow \quad x = \sqrt{c}.
\]

Then

\[
    y = \frac{c}{x} = \sqrt{c}.
\]

So the minimal sum occurs when \( x = y = \sqrt{c} \), and the minimum value of the sum is

\[
    x + y = 2\sqrt{c}.
\]

Let \( c = 25 \). Then the two positive factors that minimize the sum are

\[
    x = y = \sqrt{25} = 5,
\]

and the minimal sum is

\[
    x + y = 5 + 5 = 10.
\]

\section{Power Rule Derivation}

We derive the \emph{power rule}

\[
    \frac{d}{dx} x^n = nx^{n - 1}
\]

for \(n \in \N\) using the limit definition:

\[
    \frac{d}{dx} x^n = \lim_{h \to 0} \frac{{(x + h)}^n - x^n}{h}.
\]

Using the Binomial Theorem,

\[
    {(x + h)}^n = x^n + nx^{n-1}h + \cdots + h^n
    \Rightarrow \text{only the linear term survives after dividing by } h.
\]

\section{Implicit Differentiation}

Given

\[
    9x^2 + 4y^2 = 25,
\]

differentiate both sides implicitly:

\[
    18x + 8y \frac{dy}{dx} = 0
    \Rightarrow \frac{dy}{dx} = -\frac{18x}{8y} = -\frac{9x}{4y}.
\]

This method is used when \(y\) is not isolated and is defined implicitly.

\section{Hermit Polynomials}

A \emph{Hermit polynomial} is defined as:

\[
    H_n(x) = (-1)^n e^{x^2} \frac{d^n}{d x^n} e^{-x^2}.
\]

They appear when taking the multiple derivatives of \(e^{-x^2}\).
